\subsection{The density-level comparisons}
\label{rlevels:endpoints}

We apply both comparisons to Definition~\ref{def:r-streaming-model}.
Assume that the original uniform-prior experiment has angular success
probability at least \(2/3\), that \(0<\epsilon\le1/10\), and that
\(M=M(d)=o(d^2)\); \(T\) is its deterministic finite horizon.  Put
\(L=\log(1/\epsilon)\ge\log10\).  A worst-signal success premise of
\(2/3\) first implies this uniform-prior premise by integration.

Apply Lemma~\ref{lem:r-seed-reduction} to the original jointly
measurable experiment with \(p_0=2/3\).  Fix a seed \(\xi\) for which
\eqref{eq:r-seed-factorization} holds and conditional success is at
least \(2/3\), retain \(\mathcal L(U_0\mid\Xi=\xi)\), and then apply
Lemma~\ref{lem:r-borel-kernels}.  These reductions preserve the fixed-seed
uniform-prior joint law, the state bound, and the horizon.  All learner
statements in this subsection concern that experiment.  We prove
\(T\ge c dL\) for an absolute \(c>0\) in all sufficiently large
dimensions; the eventual threshold may depend on the memory sequence.

For completeness, \eqref{rlevels:binary-kl} applied to success
under the product signal/output law and Lemma~\ref{lem:r-caps}
give the unconditional bound
\begin{equation}\label{rlevels:output-information}
 I(S;\widehat S)\ge\tfrac23(d-1)L-\log2.
\end{equation}

Restrict to \(T<dL\).  Lemma~\ref{lem:r-stopping-reduction}
provides absorbing padded states, a state-only terminal output,
data-independent initial state \(\overline U_0\), and fresh future
blocks.  Its clock bound implies, in particular,
\begin{equation}\label{rlevels:clock}
 L\le C\frac{M+1+\log d}{d},\qquad T=O(M+d),\qquad
 H_*:=M\log2+\log(T+2)=o(d^2).
\end{equation}
The lemma also gives
\begin{equation}\label{rlevels:linear}
 T>d/4.
\end{equation}

Both comparisons below use this padded learner.  Their auxiliary
matrices and block lengths remain distinct.

For the one-level comparison, return to
\eqref{rlevels:one-parameters} and use
\(B_{\mathrm{one}}=\lceil T/k\rceil\) padded blocks.  In a block let \(Z\)
and \(U\) be the entering and exiting states.  The next matrix
\(X\) is independent of the pair \((S,Z)\).  Draw \(B\) independently
of \((S,Z,X,U)\), and put \(G=(X;B)\) as in Proposition~\ref{rlevels:one-comparison}.
Conditional on \((G,GS)\), the entire block data \((X,XS)\)
are known.  The transition kernel from \(Z\) to \(U\) then has no
additional dependence on \(S\).  Conditional data processing gives
\[
 \widetilde{\mathcal D}(U)
 \le I_{\mathrm{act}}(S;Z\mid G,GS)=\mathcal D(Z),
\]
where the equality uses independence of \(G\) from \((S,Z)\).
By \eqref{rlevels:clock}, \(H(U)\le H_*\le d^2\) eventually,
so Proposition~\ref{rlevels:one-comparison} costs at most \(Cd\)
per block.  The initial boundary \(\overline U_0\) is independent of
\(S\), and its auxiliary \(G\) is independent of the pair
\((S,\overline U_0)\); hence \(\mathcal D(\overline U_0)=0\), and
\begin{equation}\label{rlevels:one-recurrence}
 \mathcal D(U_f)\le CdB_{\mathrm{one}}
\end{equation}
for the final padded state.

Take the independent \(m\)-row matrix in the definition of
\(\mathcal D\).  Its residual radius \(R\) satisfies
\[
 \Pr\{R<1/2\}\le \frac{m/d}{3/4}\le\frac13.
\]
Lemma~\ref{lem:r-residual-terminal}, with
\(p=2/3\), \(\rho=1/2\), and \(\eta=1/3\), now gives
\begin{equation}\label{rlevels:one-endpoint}
 \mathcal D(U_f)
 \ge\frac{d-m-1}{3}(L-\log4)-\log2
 \ge c_0dL
\end{equation}
for an absolute \(c_0>0\) in sufficiently large dimension.
The final inequality is uniform for \(L\ge\log10>\log4\).
Since \(T>d/4\ge k\), the common raw-data bound makes
\(\lceil T/k\rceil\le2T/k\).  Hence
\eqref{rlevels:one-recurrence} is \(O(T)\), and comparison with
\eqref{rlevels:one-endpoint} proves \(T\ge cdL\).

For the two-label comparison, use
\eqref{rlevels:two-parameters} and
\(B_{\mathrm{two}}=\lceil T/n\rceil\) padded batches.  Let \(V_i\) be the
state after \(i\) batches.  At a batch the incoming \(A\) is
independent of \((S,V_i)\), and \(B\) is independent of the
whole actual experiment.  Apply Proposition~\ref{rlevels:two-comparison}
to the outgoing state \(V_{i+1}\); the required entropy bound follows
from \eqref{rlevels:clock} on the branch \(T<dL\).
Given \((D,DS)\), the exact batch \((A,AS)\) is known, so
the update is conditionally a channel from \(V_i\) without further
access to \(S\).  Therefore
\[
 \mathcal J(V_{i+1})
 \le I_{\mathrm{act}}(S;V_{i+1}\mid D,DS)+Cd
 \le I_{\mathrm{act}}(S;V_i\mid D,DS)+Cd
 =\mathcal J(V_i)+Cd.
\]
The last equality uses independence of \(D\) from \((S,V_i)\).
Set \(V_0=\overline U_0\).  It is independent of \(S\), and its
auxiliary \(G\) is independent of the pair \((S,V_0)\); hence
\(\mathcal J(V_0)=0\).  Iteration gives
\begin{equation}\label{rlevels:two-recurrence}
 \mathcal J(V_{B_{\mathrm{two}}})\le CdB_{\mathrm{two}}\le C'(T+d).
\end{equation}

For the independent \(k\)-row matrix defining \(\mathcal J\),
\eqref{eq:r-residual-markov} at \(\rho=1/4\) gives
\[
 \Pr\{R<1/4\}\le\frac{16}{15}\frac{k}{d}\le\frac8{15}.
\]
Apply Lemma~\ref{lem:r-residual-terminal} with \(p=2/3\) and
\(\eta=8/15\).  It gives
\begin{equation}\label{rlevels:two-endpoint}
 \mathcal J(V_{B_{\mathrm{two}}})
 \ge\frac2{15}(d-k-1)(L-\log8)-\log2.
\end{equation}
This is at least \(cdL\) uniformly for \(L\ge\log10\), since
\(L-\log8\ge(1-\log8/\log10)L>0\) and
\(d-k-1\ge d/3\).
The upper bound
\eqref{rlevels:two-recurrence} is \(O(T)\) by
\(T>d/4\).  This proves \(T\ge cdL\); the discarded branch
\(T\ge dL\) already has the same order.

In both arguments the auxiliary rows are analysis variables.  The
learner receives only its prescribed fresh Gaussian stream with exact
labels, and the comparisons permit arbitrary computation within a
transition subject to the persistent finite-state restriction.
